\section*{Chapitre \ref{ch:b2:mass-spec-atomic} --- Spectrométrie de masse et spectroscopie atomique}

\begin{solution}{exo:b2:mass-spec-atomic:1}
78~; $400/2 = 200$~; $194 + 1 = 195$, $[\mathrm M + \mathrm H]^+$.
\end{solution}

\begin{solution}{exo:b2:mass-spec-atomic:2}
Masse impaire~: un nombre impair d'atomes d'azote, un seul pour une molécule aussi petite. \ce{C3H7NO}
(propanamide, 73) ou \ce{C4H11N} (butan-1-amine, 73).
\end{solution}

\begin{solution}{exo:b2:mass-spec-atomic:3}
Hauteurs presque égales~: un brome. 3 : 1~: un chlore.
\end{solution}

\begin{solution}{exo:b2:mass-spec-atomic:4}
Jaune~: sodium, \qty{589.0}{nm} et \qty{589.6}{nm}. Rouge carmin~: lithium, \qty{670.8}{nm}.
% ledger: asd:NaD2, asd:NaD1, asd:Li671
\end{solution}

\begin{solution}{exo:b2:mass-spec-atomic:5}
$n_C \approx 6.7/1.11 \approx 6$.
\end{solution}

\begin{solution}{exo:b2:mass-spec-atomic:6}
\ce{CO}~: 27.9949~; \ce{N2}~: 28.0061~; \ce{C2H4}~: 28.0313 u. Pouvoirs de résolution~: \ce{CO}/\ce{N2} $28/0.0112
\approx 2.5 \times 10^3$~; \ce{N2}/\ce{C2H4} $28/0.0252 \approx 1.1 \times 10^3$~; \ce{CO}/\ce{C2H4}
$28/0.0364 \approx 7.7 \times 10^2$.
% ledger: mass:C12, mass:O16, mass:N14, mass:H1
\end{solution}

\begin{solution}{exo:b2:mass-spec-atomic:7}
86~: l'\omterm{def:b2:mass-spec-atomic:molecular-ion}{ion moléculaire} de \ce{C5H10O}. 57~: \ce{C2H5CO+}, \omterm{def:b2:aromatic-substitution:friedel-crafts}{ion acylium} issu d'une coupure $\alpha$ (perte
d'un radical éthyle, 29). 29~: \ce{C2H5+}.
% ledger: ms:pentan-3-one
\end{solution}

\begin{solution}{exo:b2:mass-spec-atomic:8}
86~: $\mathrm M^{+\bullet}$. 71~: $\mathrm M - \ce{CH3}$, l'acylium \ce{C3H7CO+}. 43~: \ce{CH3CO+}
(perte du radical propyle). 58~: \omterm{def:b2:mass-spec-atomic:fragmentation}{réarrangement de McLafferty}, par un hydrogène du carbone $\gamma$ de la
chaîne propyle,
avec perte d'éthène. Dans la pentan-3-one, les groupes éthyle n'ont pas de carbone $\gamma$~: pas d'ion
de McLafferty~; son petit pic à 58 est le pic isotopique \ce{^{13}C} de l'ion à 57 (trois carbones,
environ 3.3~\%).
% ledger: ms:pentan-2-one, ms:pentan-3-one
\end{solution}

\begin{solution}{exo:b2:mass-spec-atomic:9}
Une droite passant par l'origine~: pente $\sum c_iA_i/\sum c_i^2 = 1.553/30.0 \approx \qty{0.0518}{L/mg}$.
Échantillon~:
$0.128/0.0518 \approx \qty{2.47}{mg/L}$ de cuivre.
\end{solution}

\begin{solution}{exo:b2:mass-spec-atomic:10}
$t = L\sqrt{m/(2eU)} = 1.00\sqrt{1000 \times 1.6605 \times 10^{-27}/(2 \times 1.6022 \times 10^{-19}
\times 2.00 \times 10^4)} \approx \qty{16.1}{\micro s}$. Pour 1001 u, $t$ est multiplié par
$\sqrt{1.001} \approx 1 + 0.0005$~: $\Delta t \approx \qty{8.0}{ns}$.
% ledger: const:u, const:e
\end{solution}

\begin{solution}{exo:b2:mass-spec-atomic:11}
Deux chlores~: $(0.758 + 0.242x)^2$ donne $0.575$, $0.367$, $0.0586$~: $m/z$ 84, 86 et 88 dans le rapport
100 : 64 : 10.
\end{solution}

\begin{solution}{exo:b2:mass-spec-atomic:12}
\ce{^{40}Ar^{16}O+}~: $39.9624 + 15.9949 = 55.9573$~; \ce{^{56}Fe}~: 55.9349~; $R = 56/0.0224 \approx 2.5 \times
10^3$. Sinon, mesurer un autre isotope du fer, \ce{^{57}Fe}, ou éliminer \ce{ArO+} dans une cellule de
collision avant l'analyseur.
% ledger: mass:Ar40, mass:Fe56, mass:O16
\end{solution}

\begin{solution}{pb:b2:mass-spec-atomic:1}
\textbf{1.} 156, 158 et 160 (avec 157 et 159)~: la molécule existe sous plusieurs formes isotopiques.
\textbf{2.} Trois pics distants de deux unités, celui du milieu le plus grand~: un chlore et un brome.
\textbf{3.} 156.
\textbf{4.} Masse paire~: pas d'azote, ou un nombre pair.
\textbf{5.} $156 - 79 - 35 = 42$~: \ce{C3H6}.
\textbf{6.} \ce{C3H6BrCl}~; nombre d'insaturations $(2 \times 3 + 2 - 6 - 2)/2 = 0$~: ni cycle ni double
liaison.
\textbf{7.} $0.5/14.5 \approx 3.4~\%$, $3.4/1.11 \approx 3$ carbones.
\textbf{8.} Le pic à 157 est donné à 0.1 près sur une échelle où il vaut 0.5~: une incertitude de 20~\%.
\textbf{9.} \ce{C3H6^{81}Br^{35}Cl} et \ce{C3H6^{79}Br^{37}Cl}.
\textbf{10.} $3 \times 12 + 6 \times 1.007825 + 78.918338 + 34.968853 \approx 155.9341$.
\textbf{11.} $156/(156.151 - 155.934) \approx 7.2 \times 10^2$.
\textbf{12.} Oui~: deux hydrogènes de moins, 154.
\textbf{13.} Un atome de brome (79 ou 81) est perdu~: \ce{C3H6Cl+}, 77 avec \ce{^{35}Cl}, 79 avec \ce{^{37}Cl},
dans le rapport 3 : 1.
\textbf{14.} La perte de HBr~: le radical-cation $\mathrm{C_3H_5Cl^{+\bullet}}$, de masse paire.
\textbf{15.} \ce{C3H5+}, le cation allyle, après la perte des deux halogènes (Br, puis HCl).
\textbf{16.} \ce{CH2Cl+} (49 et 51, 3 : 1).
\textbf{17.} \ce{CH2Br+} (93 et 95) et \ce{C2H4Br+} (107 et 109), chacun 1 : 1~: le brome est à une extrémité de
la chaîne, le second ion provenant de la perte de \ce{CH2Cl}.
\textbf{18.} La liaison \ce{C-Br} est plus faible que la liaison \ce{C-Cl}, et \ce{Br} est le meilleur radical
partant.
\textbf{19.} Non~: pas de pics à 127--131~; \ce{CH2Cl+} et \ce{CH2Br+} montrent que chaque halogène est porté
par un \ce{CH2} terminal.
\textbf{20.} \ce{BrCH2CH2CH2Cl}, le 1-bromo-3-chloropropane.
\textbf{21.} 156/158/160~: M~; 77/79~: M $-$ Br~; 76~: M $-$ HBr~; 107/109~: M $-$ \ce{CH2Cl}~;
93/95~: \ce{CH2Br+}~; 49/51~: \ce{CH2Cl+}~; 41~: \ce{C3H5+}~; 39, 27~: \ce{C3H3+}, \ce{C2H3+}.
\textbf{22.} Non~: il n'y a pas de groupe carbonyle.
\textbf{23.} Le 1-bromo-2-chloropropane, \ce{CH3CHClCH2Br}~: la perte de \ce{CH2Br} donne \ce{CH3CHCl+} à 63/65.
\textbf{24.} \ce{C3H6^{81}Br^{37}Cl}.
\textbf{25.} $M$~: $0.758 \times 0.5069 = 0.384$~; $M+2$~: $0.758 \times 0.4931 + 0.242 \times 0.5069 =
0.496$~; $M+4$~: $0.242 \times 0.4931 = 0.119$. Rapport $\boldsymbol{\approx 3.2 : 4.2 : 1}$, environ 3 : 4 :
1~; le spectre donne $14.5 : 18.6 : 4.4 = 3.3 : 4.2 : 1$.
\end{solution}
